%-------------------------------------------------------------------------------
% N. Personal Project — VocaPatch: 한국어 어휘 학습 iOS 앱 (TestFlight 배포)
%   p.1 = 배경·목표 + 메인 스크린샷 + 시스템 구성
%   p.2 = 설계 과정 및 수행 내용
%   ※ 섹션 번호(아래 "4.")는 배치 순서에 맞게 수정
%   ※ 이미지 2장 필요: p12.jpg(앱 스크린샷 콜라주), p12pipe.png(시스템 구성 다이어그램)
%-------------------------------------------------------------------------------
\cvsection{5. Personal Project : VocaPatch — 한국어 어휘 학습 iOS 앱 (TestFlight 배포)}

\projsummary{국립국어원 사전 데이터 기반 어휘 퀴즈에 3D 캐릭터 수집·성장을 결합한 학습 앱 --- 서버 권위 경제 설계와 TestFlight 배포·운영까지}
\projfigL[7.6cm]{p12.jpg}

\projpart{\faBookOpen\hspace{1.3mm}배경 및 목표}
\begin{cvitems}
  \item {\textbf{배경}\enskip 어휘 암기 앱은 며칠 만에 이탈하는 경우가 대부분 --- 학습을 지속시키는 보상 구조가 없다는 것이 근본 문제.}
  \item {\textbf{목표}\enskip 게임의 수집·성장 시스템(3D 캐릭터·가챠·장비)을 학습 보상으로 연결한 iOS 어휘 학습 앱을 개발하고, TestFlight 배포·운영까지 완결.}
  \item {\textbf{개발 기간}\enskip 2026.05 -- 2026.06 (1인 개발, AI 코딩 도구 활용 / TestFlight 1.0.1 운영 중)}
  \item {\textbf{기술 스택}\enskip \pill{SwiftUI} \pill{Supabase} \pill{PostgreSQL} \pill{Edge Functions} \pill{국립국어원 Open API} \pill{Python} \pill{TestFlight}}
  \item {\textbf{콘텐츠 규모}\enskip 자체 구축 어휘 문항 1,050개 / 3D 캐릭터 21종 / DB 마이그레이션 무중단 운영 반영}
\end{cvitems}

\projpart{\faScrewdriverWrench\hspace{1.3mm}시스템 구성}
\projfigL[7.0cm]{p12pipe.png}

\newpage
\projpart{\faMagnifyingGlass\hspace{1.3mm}설계 과정 및 수행 내용}

\stephead{1}{국립국어원 사전 기반 콘텐츠 파이프라인}
\begin{substeps}
  \item 퀴즈 앱의 품질은 결국 콘텐츠가 결정 --- 뜻풀이·예문의 출처 신뢰성을 확보하기 위해 표준국어대사전(stdict)·한국어기초사전(krdict) Open API를 연동.
  \item \textbf{파서 설계} --- 두 API가 같은 항목을 서로 다른 XML 스키마로 반환하는 문제(필드 표기·옵션 파라미터 차이)를 흡수하는 통합 파서를 구현.
  \item \textbf{API 프록시} --- 인증 키가 클라이언트에 노출되지 않도록 사전 검색을 Supabase Edge Function 뒤로 감춤.
  \item \textbf{편집 워크플로우} --- 1,050문항 문제 은행을 스프레드시트 export/upload 스크립트로 관리해, 콘텐츠 수정을 앱 배포와 분리.
\end{substeps}

\stephead{2}{서버 권위(Server-authoritative) 경제 시스템}
\begin{substeps}
  \item XP·골드·보상을 클라이언트가 계산하면 조작과 경합(같은 보상 이중 지급)에 취약 --- 판정 권위를 전부 서버로 옮기는 설계가 필요.
  \item \textbf{원자성·멱등성} --- XP 적립과 보상 수령을 PostgreSQL RPC로 원자화하고, 세션 단위 멱등 처리와 보상 상태의 단조성(monotonic)을 DB 레벨에서 강제 --- 같은 요청이 중복 도착해도 이중 지급이 발생하지 않음.
  \item \textbf{낙관적 업데이트} --- 클라이언트는 즉시 반응하고 서버 판정과 어긋나면 롤백. 로컬 캐시(stale)가 서버 검증 액션을 가로막던 장착 버그를 클라이언트 사전 게이트 제거로 해결 --- ``서버가 검증하는 액션 앞에 클라 게이트를 두지 않는다''는 원칙을 정립.
\end{substeps}

\stephead{3}{3D 캐릭터 컬렉션·가챠 시스템}
\begin{substeps}
  \item 3D 에셋 21종을 앱에서 실시간 렌더링하면 성능·용량 부담이 큼 --- Python 에셋 파이프라인으로 GLB를 썸네일로 사전 렌더링해 앱은 이미지만 소비.
  \item \textbf{에셋 결함 복구} --- 모델 왼쪽 절반이 사라지는 미러링 누락 결함을 정점 X좌표 분포 진단 스크립트로 찾아 일괄 복구.
  \item \textbf{가챠·장비} --- 럭키박스 확률 지급과 장비 장착·효과(골드/경험치 배율)를 서버 경제 시스템과 연동해, 수집이 학습 보상으로 순환되는 구조를 완성.
\end{substeps}

\stephead{4}{학습 게이미피케이션 UX와 배포 운영}
\begin{substeps}
  \item \textbf{실시간 피드백} --- XP·콤보·스트릭을 즉각 시각화. 연속 정답 시 보상 애니메이션이 서로를 덮어쓰는 타이머 경합을 상태 seed 검사 패턴으로 방지.
  \item \textbf{통계 무결성} --- '다시 풀기' 세션을 플래그로 격리해 XP·콤보·정답률 통계를 오염시키지 않도록 설계 --- 재학습을 장려하면서도 기록의 신뢰를 유지.
  \item \textbf{배포·운영} --- TestFlight 1.0.1 배포, 스키마 변경을 마이그레이션 스크립트로 무중단 PROD 반영.
\end{substeps}
